\begin{ejercicio}
    % ID: 41
    Mostrar que la inclusión $\iota:S^n\hookrightarrow \R^{n+1}-\{0\}$ pasa al los cocientes y define un homeomorfismo canónico $\overline{\iota}:S^n/\Z_2\to \R P^n$ tal que el siguiente diagrama conmuta
    \begin{equation*}
        \begin{tikzcd}[column sep=1cm, row sep=1cm]
            S^n\arrow[r,"\iota", hook]\arrow[d,"\pi_1"] & \R^{n+1}-\{0\}\arrow["\pi_2",d]\\
            S^n/\Z_2\arrow["\overline{\iota}",r] & P \R^n
        \end{tikzcd}
    \end{equation*}
    y que $S^1/\Z_2\cong S^1,$ por lo tanto, $\R P^1\cong S^1.$
\end{ejercicio}

\begin{demostracion}

\end{demostracion}
